\documentclass[11pt]{article}
\usepackage[margin=1in]{geometry}
\usepackage{amsmath,hyperref}
\begin{document}
\begin{center}\large Prescribed chemistry in Ciais et al. (2026)\end{center}
\noindent\textbf{Reference.} Ciais et al., \emph{Science} 391, eadx8262,
\href{https://doi.org/10.1126/science.adx8262}{doi:10.1126/science.adx8262}.
The full article and 46-page supplement (Figs.~S1--S18, Table~S1) were read.
Numerical comparisons retain 2019--2022.

\paragraph{Method.} Three GOSAT inversions prescribe two independently
constructed OH fields. OH-1 combines observed precursor concentrations with
grid-cell chemical-box sensitivities evaluated for 2005. OH-2 uses ERA5-driven
LMDz--INCA chemistry with updated precursor emissions, including aviation;
its magnitude is multiplied by 0.88 to match the earlier TRANSCOM reference.
OH-1 assumes fixed NO$_x$ vertical profiles. PYVAR--LMDz optimizes weekly total
fluxes at $1.9^\circ\times3.75^\circ$; WESTLAKE optimizes seasonal land fluxes
at $4^\circ\times5^\circ$; GONGGA optimizes four source categories at
$2^\circ\times2.5^\circ$ in 14-day windows, with an initial-state adjustment
using the first seven days. Wetland priors have no interannual variation.
Quoted spreads combine three models and two OH fields, without sampling
all uncertainty in the prescribed chemistry.

\paragraph{Findings.} Emissions increase by $22.0\pm3.5$ Tg CH$_4$ yr$^{-1}$
in 2020 relative to 2019. Wetland and inland water emissions average
$8.6\pm2.6$ Tg CH$_4$ yr$^{-1}$ above 2019 in 2020--2022. Africa and
Southeast Asia dominate the increase; southern South American wetlands
instead decrease by $1.4\pm0.8$ Tg CH$_4$ yr$^{-1}$.
The headline 83\% OH contribution uses 2019--2023, outside our target interval.
The supplement defines it as the mean of four fractions of adjacent-year
growth changes explained by changes in OH removal. It is neither a fraction
of cumulative growth nor a fixed-2019 forcing fraction. Separate reported
OH--growth regression $R^2$ values are 0.5--0.7.

\paragraph{Sector interpretation.} In mixed grid cells, emission changes are
reallocated when the agricultural share differs by more than 20\% from the
corresponding bottom-up fraction. Thus wetland versus agricultural
attribution is partly prescribed. Prior-proportional allocation would assign
$46.8\pm14.8$\% of global changes to agriculture/waste and
$34.7\pm4.6$\% to wetlands. A global isotope-box calculation favors GONGGA's
source mixture, conditional on source signatures and fractionation assumptions.
The supplement inconsistently describes GONGGA's treatment of other natural
sources on pp.~4--5; those components are not independently reconstructed here.

\paragraph{Data and comparison.} Table~S1 supplies bottom-up baseline levels
and period-mean anomalies. Annual inversion emissions and OH-removal anomalies
are available graphically. We digitize central values and separate reading
precision from ensemble spread. Native annual burdens and absolute 2019 OH
removal are not tabulated. Conversion to He's OH forcing therefore uses an
explicit common-burden/lifetime scenario. The air-mass-weighted OH maps cannot
replace reaction-weighted loss frequencies. No inversion is repeated.
\end{document}
